% Generated draft: requires mathematical and source review.
\documentclass[11pt]{article}
\usepackage[margin=25mm]{geometry}
\usepackage{fontspec}
\setmainfont{Latin Modern Roman}
\usepackage{amsmath,amssymb,mathtools,unicode-math}
\setmathfont{Latin Modern Math}
\usepackage{enumitem,needspace,xurl,hyperref}
\setlength{\parindent}{0pt}
\setlength{\parskip}{5pt}
\setlength{\emergencystretch}{3em}
\newcommand{\N}{\mathbb{N}}
\newcommand{\Z}{\mathbb{Z}}
\newcommand{\Q}{\mathbb{Q}}
\newcommand{\R}{\mathbb{R}}
\newcommand{\C}{\mathbb{C}}
\newcommand{\F}{\mathbb{F}}
\newcommand{\A}{\mathbb{A}}
\newcommand{\PP}{\mathbb{P}}
\newcommand{\E}{\mathbb{E}}
\newcommand{\eps}{\varepsilon}
\newcommand{\dd}{\,\mathrm{d}}
\begin{document}
\section*{UnsolvedMath ID: 10400076}
\textbf{Source problem code:} catalogue entry
\subsection{Definitions and mathematical statement}
The primary request is to determine the image and kernel of an ``above map'' \(\tau\), whose definition is omitted. The summary describes it schematically as a map from \(\mathcal A_\Sigma\) to functions on \(\mathcal M_G(\Sigma)\). The diagram algebra, group, moduli space, function category and construction are absent, so they remain placeholders. Once these are fixed, the target is a complete description of
\[
 \operatorname{im}\tau=\{\tau(a):a\in\mathcal A_\Sigma\},\qquad
 \ker\tau=\{a\in\mathcal A_\Sigma:\tau(a)=0\}
\]
in the intended linear/algebraic setting, not merely examples of values or relations.

The excerpt supplies conformal-block background. Let \(p_1,\ldots,p_n\) be marked points on the intended Riemann surface \(\Sigma\), let \(\mathfrak g\) be the Lie algebra and \(H_i\) representations of its affine extension \(\widehat{\mathfrak g}\). Let \(\mathfrak g_{\mathrm{out}}\) consist of meromorphic \(\mathfrak g\)-valued functions with poles only at the marked points, acting on \(V=H_1\otimes\cdots\otimes H_n\) by Laurent expansions. The given conformal blocks are
\[
 \{\phi:V\to\mathbb C\text{ linear}:\phi(a\cdot v)=0
 \text{ for all }a\in\mathfrak g_{\mathrm{out}},v\in V\}.
\]
Thus in this formulation they are dual to the coinvariant quotient \(V/\mathfrak g_{\mathrm{out}}V\), not literally the quotient itself. The twisted version replaces meromorphic functions by meromorphic sections of a \(\mathfrak g\)-local system.

The levels, module restrictions, coordinates/local-system data and relationship of these blocks to \(\tau\) are not supplied. No kernel formula is inferred merely from this background. The summary says image/kernel are partially understood but gives no exact theorem; that partial-progress label remains an unverified source assessment pending primary-literature review.
\subsection{Short English statement}
Determine the full image and kernel of the source's map \(\tau\). The attached WZW conformal-block definition is retained, but the map and its connection to that background are missing.
\subsection{Sources}
\begin{itemize}
\item Local registry: \url{data/unsolvedmath/problems.json}, permanent UnsolvedMath ID \(10400076\), source code catalogue entry; supplied statement, background, and research summary, updated 2026-08-06.
\item Ohtsuki, \emph{Problems on invariants of knots and 3-manifolds} (2002), Problem 3.26, attributed to T. Kohno, PDF page 72: \url{https://msp.org/gtm/2002/04/gtm-2002-04-024s.pdf}. Supplied PDF-font-structure extraction accessed 2026-07-29; provenance distinguishes 231 numbered items.
\item No equation for \(\tau\), precise partial result or separate resolution reference is supplied.
\end{itemize}
\end{document}
